\documentclass[11pt, a4paper, onecolumn]{article}

%install packages as needed here
%\usepackage{geometry}  %useful for changing page layout
%\geometry{top=2.5cm, bottom=2.5cm, left=2.5cm, right=2.5cm}

% \title{\textbf{Astrophysical simulation of N-Bodies}
% \author{\textbf{Darragh Morrissey}$^{1}$ \\
%        \small $^{1}$Queen's University, Belfast \\}}
% \date{\today}

\title{Simulation of the Motion of Astronomical Bodies}
\author{Darragh Morrissey - 40451054}

\begin{document}

\maketitle
Here's the outline for the report regarding the layout and process. 
\begin{itemize}
    \item Structured work:
    \begin{itemize}
        \item Structured(1) - Stationary Sun and moving Earth
        \item Structured(2) - Track period of Earth's orbit
        \item Structured(3) - Track period of new planet based off student id number
    \end{itemize}
    \item Allow Sun to move due to pull of Earth
    \item Extend simulation to include all planets of the Solar System
    \item Allow planets of Solar System to move clockwise and counterclockwise to view how stability is affected
    \item Add asteroids to simulation to examine bahviour of kirkwood gaps 
\end{itemize}

\section{Abstract}

Using python, a simulation of the bodies of the Solar System was created  

\section{Introduction}
The aims of this project was to accurately model the behaviour of astronomical bodies in the Solar System and to examine the nature of Kirkwood gaps using python. 
It was assumed that all motion will take place on the x-y plane 


\section{Methods}
Usage of Runge Kutta method

Sun and Earth model with stationary Sun

Orbital time tracker

Sun and Earth model with Sun able to move

Sun and arbitrary body model where initial velocity is given to body but not Sun. 

Give Sun initial momemtum to give system 0 total momentum

Add rest of bodies from Solar System into simulation

Examine Kirkwood gaps by only considering Sun, Jupiter and asteroids at various starting orbits

\section{Results}





\section{Conclusion}





\section{Appendix}





\end{document}